AI creative systems treat domain expertise as the primary persona variable:
add a puzzle-designer label, get better puzzle design.
This assumption is wrong in a specific and actionable way.
In a controlled authoring study, \textit{``you are an artist''} outperformed
\textit{``you are a puzzle designer''} by 15.6 points on a 45-point weighted rubric,
and was the only condition to escape mechanism gravity into structural novelty.
The structural explanation for this gap is unknown.

We synthesize three empirical studies across authoring (trait decomposition) and reviewing
(profile ablation) roles in puzzle-hunt design, identifying a shared three-tier quality
architecture from evidence across both roles.

A three-tier structure governs creative quality in both roles.
Tier~1 (Orientation) is irreplaceable: removing the experience-first stance causes catastrophic
authoring failure (L-T1 $= -29.4$ points) and sub-baseline reviewing performance
(C5 lens-only $= 21.7$, below bare baseline C0 $= 23.7$).
Tier~2 (Lens) is load-bearing: elegance drive and rigor gate threshold crossing in authoring;
three non-redundant principles perform the same function in reviewing.
The minimum sufficient combination (T1+T4+T5) crosses the publishable threshold at 35.7
on its first reconstruction condition (R3).
Tier~3 (Specificity) is amplifying: domain expertise and the review lens add ceiling performance
but cannot substitute for lower tiers.
The artist label activates Tier~1; the domain label activates Tier~3.
The 15.6-point gap is the empirical cost of starting at the wrong tier.

We formalize the OLE formula (Orientation $\to$ Lens $\to$ Expertise) as a minimum
sufficient creative prompt construction protocol.
Three sentences constitute the irreducible prompt for any domain:
\textit{``Think first about how the other person will experience this.
Remove anything unnecessary.
Verify every claim.''}
This is the first cross-role synthesis of creative quality architecture in AI prompting,
providing empirical grounding for a prompt design principle applicable to AI creative tasks
in domains where quality is defined by receiver experience rather than domain convention.
We predict this principle generalizes across such domains (cross-domain validation pending).
